% TOP_PROBLEM: {"id": 1469, "title": "Serre's Positivity Conjecture", "category_id": 4, "category": {"id": 4, "name": "algebra", "display_name": "Algebra", "description": "Group theory, ring theory, field theory, and algebraic structures.", "slug": "algebra", "order_index": 4, "created_at": "2026-07-31T15:26:25.671Z"}, "status": "open"}
% FIRST_POSTED: 2026-09-27
% !TeX program = lualatex
% ATTEMPT_STATUS: unresolved
\documentclass[11pt]{article}
\usepackage[margin=25mm]{geometry}
\usepackage{fontspec}
\setmainfont{Latin Modern Roman}
\usepackage{amsmath,amssymb,mathtools}
\usepackage{unicode-math}
\setmathfont{Latin Modern Math}
\usepackage{xurl,hyperref}
\hypersetup{colorlinks=true,urlcolor=blue}
\setcounter{secnumdepth}{0}
\setlength{\parindent}{0pt}
\setlength{\parskip}{5pt}
\setlength{\emergencystretch}{3em}
\begin{document}
\section*{\texorpdfstring{Serre's Positivity Conjecture}{Serre's Positivity Conjecture}}\label{1469}
\subsection{Definitions and mathematical statement}
Let $(R,\mathfrak m)$ be a commutative regular Noetherian local ring and let $M,N$ be nonzero finitely generated modules with $M\otimes_RN$ of finite length. Their dimensions are dimensions of their supports. Since $R$ has finite global dimension, the alternating sum
\[
 \chi^R(M,N)=\sum_{i\ge0}(-1)^i
 \operatorname{length}_R\operatorname{Tor}_i^R(M,N)
\]
is finite; the support hypothesis makes each term finite. Serre positivity asserts $\chi^R(M,N)>0$ whenever $\dim M+\dim N=\dim R$. No equal-characteristic or unramified assumption is added. Strict positivity is stronger than nonnegativity, while the complementary dimension regimes are not part of this selected statement.
\par\smallskip\textit{Formulation and concept references:} \href{https://arxiv.org/html/2404.17572}{[S1]}.

\subsection{Short English statement}
When two modules over a regular local ring meet in finite length and their dimensions add correctly, must their intersection multiplicity be strictly positive?
\subsection{Research attempt: exact Koszul positivity and the prime-module gap}
\textbf{Status.} Strict positivity for arbitrary complementary-dimensional modules over an arbitrary regular local ring is \textbf{UNRESOLVED} by this attempt. We prove several cases by explicit resolutions, including genuinely ramified examples, and give a precise reduction to prime quotients using an explicitly identified known vanishing theorem. The unresolved step is not removed by replacing the Euler characteristic with the length of the ordinary tensor product.

\textbf{Source audit.} The inspected source [S1] resolves to arXiv:2404.17572v2, dated 5 November 2024. Its Section 1.1 distinguishes nonnegativity and vanishing, known in general, from strict positivity in the ramified case. Theorem A establishes a new positivity case subject to Serre liftability; it does not assert that every module admits such a lift. We retain that hypothesis when discussing the result. The calculations here are independent elementary subcases and are not claimed to improve that theorem. In particular, no equal-characteristic argument is silently transferred to a ramified ring.

\subsubsection{1. Why all the lengths exist and why the signs matter}
Put $d=\dim R$. For finite modules $M,N$, localization commutes with Tor. If a prime $\mathfrak p$ does not lie in both supports, one of $M_\mathfrak p,N_\mathfrak p$ vanishes, and so does every localized Tor. Thus
\[
 \operatorname{Supp}\operatorname{Tor}_i^R(M,N)
 \subseteq\operatorname{Supp}M\cap\operatorname{Supp}N
 =\operatorname{Supp}(M\otimes_RN).
\]
The last equality for finite modules follows locally from Nakayama: over a local ring the tensor product of two nonzero finite modules has nonzero reduction modulo the maximal ideal. The finite-length assumption therefore puts all Tor supports at $\mathfrak m$. Finite generation then gives finite lengths. Regularity supplies finite free resolutions of length at most $d$, so the alternating sum is finite.

If $F_\bullet\to M$ is a free resolution, then $\chi^R(M,N)$ is the Euler characteristic of the homology of $F_\bullet\otimes N$. The terms of this complex usually have infinite length, so one may not simply alternate their lengths. The familiar cancellation identity for finite-dimensional complexes applies to those terms only when they themselves have finite length.

For a short exact sequence in either module variable, the Tor long exact sequence gives additivity of $\chi$, provided all relevant tensor products have finite length. To prove this, truncate the long exact sequence after the finite global-dimension bound and use the alternating length identity for an exact sequence of finite-length modules. This additivity will be used carefully below.

\subsubsection{2. The finite-length versus full-dimensional case}
Suppose $M$ has finite length and $\dim N=d$. A regular local ring is a domain, so $N$ has positive rank $r$: full-dimensional support contains the zero prime. Take a finite free resolution $F_\bullet\to N$. Now every term $M\otimes F_i$ has finite length, equal to $\operatorname{rank}(F_i)\ell(M)$. Euler characteristic of terms equals that of homology, yielding
\[
 \chi^R(M,N)=\ell(M)\sum_i(-1)^i\operatorname{rank}(F_i).
\]
Tensoring the resolution with the fraction field shows that the sum is $r$. Therefore
\[
 \chi^R(M,N)=r\,\ell(M)>0.                             \tag{1}
\]
This includes mixed characteristic and ramification. Symmetry of Tor gives the same conclusion with $M,N$ exchanged. If $d=0$, $R$ is a field and the assertion reduces to the positive product of two vector-space dimensions.

For $d=1$, $R$ is a discrete valuation ring. Complementary dimensions force one module to have dimension zero and the other dimension one, so (1) proves the entire selected statement in dimension one. The calculation is also visible from the structure theorem
$N\cong R^r\oplus\bigoplus_jR/(\pi^{a_j})$: finite torsion summands contribute equally to Tor degree zero and degree one against a finite-length module, hence contribute zero to $\chi$.

This does not extend by simply assigning a rank to lower-dimensional modules: their rank over the fraction field is zero, while proper intersections between them can have positive multiplicity. The useful ranks then occur along their support primes, which is one reason the prime-filtration reduction is natural.

\subsubsection{3. A Koszul criterion for strict positivity}
Suppose $x_1,\ldots,x_c$ is an $R$-regular sequence, set
$M=R/(x_1,\ldots,x_c)$, and assume $N$ is a nonzero Cohen--Macaulay module of dimension $c$ such that $M\otimes N$ has finite length. The finite quotient condition says that $(x_1,\ldots,x_c)$ is an ideal of definition on $N$, so it is a system of parameters. The standard Cohen--Macaulay parameter theorem says that every system of parameters of a Cohen--Macaulay module is a regular sequence on that module. This theorem is the explicitly used commutative-algebra input.

The Koszul complex $K(x;R)$ is a free resolution of $M$. Tensoring with $N$ gives $K(x;N)$. Because $x$ is $N$-regular, its higher homology vanishes. Hence
\[
 \operatorname{Tor}_i^R(M,N)=0\ (i>0),\qquad
 \chi^R(M,N)=\ell\bigl(N/(x_1,\ldots,x_c)N\bigr)>0.     \tag{2}
\]
The quotient is nonzero by Nakayama because all $x_i$ lie in $\mathfrak m$ and $N\ne0$. Since $R$ is Cohen--Macaulay, $\dim M=d-c$, so the dimension equality required in the catalogue is satisfied.

For clarity, the exactness assertion for a regular sequence can be proved inductively. For one element, the complex is $0\to N\xrightarrow{x_1}N\to0$, with injective first map. For several elements, tensor the preceding Koszul complex with the two-term complex for $x_c$. The mapping-cone long exact sequence identifies positive homology with the kernel of $x_c$ on the previous degree-zero quotient, which is zero by regularity. Degree zero is the new quotient. This proves the homological step in (2) without a hidden Tor-independence assumption.

The same argument works whenever the defining $R$-regular sequence is known separately to be $N$-regular; Cohen--Macaulayness is one sufficient way of ensuring it. We do not assert that every prime ideal is a complete intersection, or that every module in the problem is Cohen--Macaulay.

\subsubsection{4. Concrete ramified examples and exact multiplicities}
Fix a prime $p$ and integers $a,b\ge1$. Consider the complete local ring
\[
 R=\mathbb Z_p[[x,y]]/(p-x^ay^b).
\]
Its maximal ideal is generated by $x,y$, and its dimension is two: the ambient regular local ring has dimension three and we quotient by the nonzero nonunit $p-x^ay^b$. Its embedding dimension is two, so it is regular. It has mixed characteristic $(0,p)$, and $p=x^ay^b\in\mathfrak m^2$, so it is ramified. For the characteristic assertion, regularity makes the quotient a domain, and the independent parameter classes $x,y$ are nonzero. Thus $p=x^ay^b$ is nonzero, so the map from $\mathbb Z_p$ is injective.

The elements $x,y$ form a regular system of parameters, hence a regular sequence in either order. For positive $u,v$ take
\[
 M=R/(x^u),\qquad N=R/(y^v).
\]
They both have dimension one, and $x^u$ is a nonzerodivisor on $N$, because powers of a regular system of parameters remain a regular sequence. Equation (2) gives
\[
 \chi^R(M,N)=\ell(R/(x^u,y^v))=uv.                    \tag{3}
\]
The length formula can be proved by successive filtrations: modulo $y^v$, multiplication by $x$ gives $u$ successive quotients isomorphic to $R/(x,y^v)$, and modulo $x$ the parameter $y$ gives $v$ residue-field quotients. No claim is made that the Artin quotient is a vector space over $\mathbb F_p$ in every choice of $a,b,u,v$; its characteristic can be a higher power of $p$. The length argument avoids that false simplification.

In equal characteristic, the analogous example over $k[[x_1,\ldots,x_d]]$ uses complementary groups of coordinate powers. The intersection multiplicity is the product of all their exponents. These elementary models confirm that ramification itself does not force a sign problem. The difficulty lies in general modules and the possible higher Tor terms, not in the existence of the residue characteristic inside $\mathfrak m^2$ alone.

\subsubsection{5. A non-Cohen--Macaulay summand exhibits genuine cancellation}
Let $R=k[[x,y]]$, $M=R/(x)$, and
$N=R/(y)\oplus k$. The module $N$ has dimension one but is not Cohen--Macaulay, because it contains the finite-length summand $k$. The resolution
$0\to R\xrightarrow{x}R\to M\to0$ gives
\[
 \operatorname{Tor}_0(M,N)\cong k\oplus k,\qquad
 \operatorname{Tor}_1(M,N)\cong k,
\]
and all higher Tor groups vanish. Indeed multiplication by $x$ is injective on $R/(y)\cong k[[x]]$ and zero on $k$. Thus
\[
 \ell(M\otimes N)=2,\qquad \chi^R(M,N)=2-1=1.          \tag{4}
\]
All catalogue hypotheses hold, and positivity survives. But ordinary tensor-product length is not the answer. Any proposed proof which throws away higher Tor terms would already miscompute this example.

The complementary dimension condition is also essential. Over the DVR $k[[t]]$, take $M=N=k$. Then Tor in degrees zero and one is $k$, so $\chi(k,k)=0$, although $M\otimes N\ne0$. Here $\dim M+\dim N=0<1$. The same cancellation is compatible with the general vanishing theorem and is not a counterexample to strict positivity at complementary dimension.

\subsubsection{6. Reduction to prime quotients, with the external theorem exposed}
The following reduction uses the known vanishing theorem for regular local rings:
\[
 \dim A+\dim B<d,\quad\ell(A\otimes B)<\infty
 \quad\Longrightarrow\quad\chi^R(A,B)=0.              \tag{5}
\]
Its full mixed-characteristic form is deep; [S1, Section 1.1] attributes it to Gillet--Soul\'e and Roberts. It is an external input here, not a consequence of our elementary Koszul calculation.

Every finite module has a finite prime filtration, whose successive quotients are $R/\mathfrak p$. Additivity applied to prime filtrations of both $M$ and $N$ expresses $\chi(M,N)$ as a sum of $\chi(R/\mathfrak p,R/\mathfrak q)$, with repetitions allowed. All relevant tensor products have finite length because their supports lie in the original supports. Moreover
$\dim R/\mathfrak p\le\dim M$ and
$\dim R/\mathfrak q\le\dim N$.

When $\dim M+\dim N=d$, every pair in which one dimension is strictly smaller vanishes by (5). Grouping the surviving terms gives
\[
 \chi^R(M,N)=
 \sum_{\substack{\dim R/\mathfrak p=\dim M\\
                  \dim R/\mathfrak q=\dim N}}
 m_{\mathfrak p}(M)m_{\mathfrak q}(N)
 \chi^R(R/\mathfrak p,R/\mathfrak q),                 \tag{6}
\]
where the coefficients are positive integers for primes actually occurring. They can equivalently be expressed as the appropriate generic lengths; the filtration form suffices for the argument.

Each nonzero module has at least one support prime of maximal dimension, so the sum has at least one pair with positive coefficient. Consequently positivity for every complementary prime-quotient pair implies the full assertion. Conversely, prime quotients are special modules, so the full assertion implies that prime-pair statement. This is an actual reduction, but it leaves arbitrary singular integral supports; it does not make them smooth complete intersections.

Nonnegativity, also known in general as discussed in [S1], would turn (6) into a sum of nonnegative quantities. It still gives no reason that at least one summand is strictly positive. Zero is not excluded by any amount of bookkeeping of signs alone.

\subsubsection{7. Lifting as a conditional route, not a universal construction}
Theorem A of [S1] gives positivity for a regular quotient $R=Q/(f)$ of an unramified regular local ring when $f\in\mathfrak m_Q\setminus\mathfrak m_Q^2$ and the relevant module admits a Serre lift to $Q$. A Serre lift is a finite $Q$-module $\widetilde M$ with $\widetilde M\otimes_QR\cong M$ and the precise dimension increase
\[
 \dim Q-\dim\widetilde M=\dim R-\dim M.
\]
This dimension equality is additional information. Merely viewing $M$ as a $Q$-module via the quotient leaves its dimension unchanged and does not provide the required lift. Nor does lifting matrices in a presentation automatically control the dimension or homology of the lifted presentation.

The source explicitly records the absence of a known counterexample to universal Serre liftability in the setting it studies, not a theorem that every module has the lift. Therefore a proof strategy based on this theorem must either construct a lift with its dimension verified or leave that construction as a gap. The explicit ramified Koszul examples above need no such unproved construction.

\subsubsection{8. Verification and the first unresolved step}
The diagnostics use exact complexes and lengths: the DVR cancellation is $(1,1)$ in Tor degrees $(0,1)$, the non-Cohen--Macaulay example is $(2,1)$, and the complementary parameter-power examples give the product in (3). Finite enumeration of lengths in these regular-sequence examples cannot test arbitrary prime ideals or all mixed-characteristic local rings.

The strongest results proved directly are the finite-length/full-rank case, all one-dimensional cases, and complete-intersection intersections with a Cohen--Macaulay complementary module, including ramified rings. With the declared vanishing theorem, the problem reduces to strict positivity for arbitrary complementary prime quotients. The first remaining step is to establish that positivity for such singular supports in the ramified setting, or supply the required lift for each of them. No argument here provides it. The general assertion is \textbf{UNRESOLVED} by this attempt.

\subsection{Sources}
\begin{itemize}
\item[S1]N. KC and A. J. Soto Levins, On liftings of modules of finite projective dimension, Section 1.1 (2024).
\url{https://arxiv.org/html/2404.17572}.
\textit{Formulation source.}
\end{itemize}
\end{document}
